% =============================================================================
% Section 3 — Data and Surface Characterization
% PINN-ONB01 Phase 1 manuscript (IJHMT submission)
% Target length: ~700 words (equations and tables excluded)
% =============================================================================

\section{Data and Surface Characterization}\label{sec:data}

Training and validating the PINN required a curated pool-boiling
corpus, which we assembled as described below. This section
documents a four-stage curation pipeline --- selection, digitization,
standardization, and per-surface bookkeeping --- that produced the
dataset summarized in \cref{tab:dataset}. We release the dataset
together with the per-figure preprocessing records to allow
independent auditing.

% -----------------------------------------------------------------------------
\subsection{Curation pipeline}\label{subsec:curation}\label{subsec:selection}\label{subsec:digitize}\label{subsec:standardize}
% -----------------------------------------------------------------------------

Candidates were drawn from a structured literature survey that
screened \num{25} primary references, of which seven
(\citealp{betz2013,bourdon2012,bourdon2015,jabardo2009,jones2009,jo2011,phan2009})
ultimately supplied data. Four criteria were applied:
(i) pool boiling on a planar or large-cylinder heater (so that a 1D
treatment is justified);
(ii) a published boiling curve from which the ONB can be localized;
(iii) at least one quantitative surface descriptor ($R_a$, $\theta$,
or an unambiguous coating identity);
(iv) an explicitly stated working fluid and pressure.
Seven primary sources passed all four criteria. They span
hydrophilic/superhydrophobic patterns~\citep{betz2013}, grafted
SAMs~\citep{bourdon2012,bourdon2015}, roughness ladders on copper,
brass, and stainless steel~\citep{jabardo2009}, EDM-textured
copper~\citep{jones2009}, hydrophilic/hydrophobic
stripes~\citep{jo2011}, and nanoparticle oxides~\citep{phan2009}. Four
fluids are represented: water, R-123, R-134a, and FC-77.

The selected boiling curves were then digitized with
WebPlotDigitizer~v5.x~\citep{rohatgi_wpd}.
Axes were calibrated from tick labels, and logarithmic scales were
verified by hand before point-by-point tracing. Mask-based
auto-detection was reserved for densely sampled monochrome curves
and always followed by visual confirmation. We required each series
to be monotonic in $q''$; violations were flagged for review rather
than deleted. The digitized numerical ranges were cross-checked
against the values reported in the source text. Raw extraction
metadata are preserved alongside each trace. The protocol produced
\num{53} raw curve files containing \num{1361} points in total
(mean \num{25.7} points per figure).

The digitized traces were converted to a common schema using an
open-source standardization script. The script converts units to SI
(e.g., kW/m$^{2} \to$ W/m$^{2}$). When the heat-transfer
coefficient $h$ is reported on either the abscissa or the ordinate,
the script reconstructs $\Delta T_{\mathrm{wall}} = q''/h$.
Per-paper figure schemas with column-by-column unit specifications
are released in the supplementary material. The standardized
entries record source paper, figure reference, surface identifier,
working fluid, wall superheat, subcooling, heat flux, roughness,
contact angle, surface category, an ONB flag, and free-text notes.

% -----------------------------------------------------------------------------
\subsection{Surface cards and ONB labels}\label{subsec:cards}\label{subsec:sfc}\label{subsec:onb_label}
% -----------------------------------------------------------------------------

Each unique heater is registered as a \emph{surface card}, and
\num{49} cards are populated. A card lists material, surface
treatment, $R_a$, $\theta$, fluid-side environment, source paper, and
any literature-reported cavity statistics ($r_c$, $N_s$). Some
surfaces are tested with multiple fluids. For example, the polished
copper substrate of \citet{jabardo2009} was tested in both R-123 and
R-134a; the corresponding surface card therefore carries a
one-to-many fluid mapping. Ten high-level categories, grouped by
source paper and material, support stratified sampling. The
category list and per-card entries are released as part of the
supplementary material.

Building on this per-surface bookkeeping, the pipeline produced
\num{82} ONB labels. Of these, \num{24} were
assigned manually from clearly identifiable inflection points (see
Fig.~4 of BETZ, Fig.~3 of JO, Fig.~12a of PHAN, Fig.~4 of
BOURDON-2012, Fig.~3 of BOURDON-2015, and Fig.~9a/9c of JABARDO).
The remaining \num{58} were assigned by a slope-change heuristic in
$\log q''$--$\log\Delta T_{\mathrm{wall}}$ space: the heuristic
selects the first point at which the local slope exceeds the
natural-convection slope by a factor of approximately \num{1.5}.
Nine heuristic labels were overridden when they fell outside the
physically admissible interval
$2 \le \Delta T_{\mathrm{ONB}} \le \SI{30}{\kelvin}$ or violated the
Hsu envelope; in each case the manual label was retained. For each
label, the procedure of \cref{subsec:hsu} infers an effective $r_c$
that is retained as a weak supervisory signal. The pipeline produced
two release tables: the full curated curve table (\num{1361} rows,
14 columns) and the ONB-label subset (\num{82} rows, 12 columns).

% -----------------------------------------------------------------------------
\subsection{Dataset statistics}\label{subsec:stats}
% -----------------------------------------------------------------------------

The curated dataset (\cref{tab:dataset}) covers \num{49} surfaces,
\num{1361} curve points, and \num{82} ONB labels across four fluids.
The incipience superheat spans
$\SI{0.7}{\kelvin} \le \Delta T_{\mathrm{ONB}} \le \SI{24.9}{\kelvin}$,
covering most of the canonical water window
(\SIrange{2}{30}{\kelvin}) and extending below it into the
low-superheat regime characteristic of refrigerants. After removing
points beyond the critical heat flux (CHF) of each surface, the
labeled heat flux spans approximately
\SIrange{0.1}{1434}{\kilo\watt\per\square\meter}. The fluid
breakdown is \num{33} water, \num{34} R-134a, \num{10} R-123, and
\num{5} FC-77. FC-77 entries are released but excluded from PINN
training because CoolProp does not provide a validated equation of
state for this fluid; the effective training count is therefore
$n=\num{77}$. Unlike purely data-driven models, the present
framework requires rigorous saturated-state properties
($\rho_{\ell}$, $\rho_v$, $h_{fg}$, $\sigma$, $c_{p,\ell}$, $k_{\ell}$)
at the operating temperature to evaluate the PDE residual, the
Hsu nucleation discriminant, and the non-dimensional scales of
\cref{subsec:nondim}. Excluding FC-77 is therefore a measure to
preserve physical rigour rather than a curation choice;
the released CSVs retain the FC-77 rows so that the exclusion
boundary remains transparent and auditable. All water experiments were conducted at
near-atmospheric pressure. The JABARDO refrigerant data span
reduced pressures $p_r \approx \numrange{0.06}{0.26}$ and provide
the only pressure variation available in the present corpus.

\begin{table}[!htbp]
\centering
\caption{Curated pool-boiling dataset by source paper. Curve points are
post-standardization counts in \texttt{boiling\_curves.csv}; ONB labels
are the rows in \texttt{onb\_dataset.csv}. The wall-superheat range is
reported over the ONB labels of each paper.}
\label{tab:dataset}
\footnotesize
\begin{tabular}{lccccc}
\hline
Source & Surfaces & Curves & ONB & Fluids & $\Delta T_{\mathrm{ONB}}$ [K] \\
\hline
BETZ\_2013      &  7 & 168 & 10 & water                  & 3.1--11.4 \\
BOURDON\_2012   &  5 & 144 &  6 & water                  & 4.5--11.0 \\
BOURDON\_2015   &  4 &  75 &  5 & water                  & 1.3--10.5 \\
JABARDO\_2009   & 20 & 589 & 44 & R-123/R-134a           & 0.7--24.9 \\
JONES\_2009     &  6 & 320 & 10 & water/FC-77            & 6.1--17.9 \\
JO\_2011        &  2 &  14 &  2 & water                  & 6.9--8.5  \\
PHAN\_2009      &  5 &  51 &  5 & water                  & 8.5--18.9 \\
\hline
Total           & 49 & \num{1361} & 82 & 4 fluids              & 0.7--24.9 \\
\hline
\end{tabular}
\end{table}
